% !TeX program = xelatex
\documentclass[12pt]{article}
\usepackage[spanish,es-nodecimaldot]{babel}

\usepackage{fontspec}
\IfFontExistsTF{Times New Roman}{\setmainfont{Times New Roman}}{\setmainfont{Liberation Serif}}

\usepackage[letterpaper,top=2.4cm,bottom=2.4cm,left=2.4cm,right=2.4cm]{geometry}
\usepackage{setspace}
% Renglon cerrado
\singlespacing
\usepackage{xurl}
\usepackage[hidelinks]{hyperref}
\usepackage{enumitem}
\usepackage{booktabs}
\usepackage{float}
\usepackage{tabularx}
\usepackage{array}
\usepackage{titlesec}

\setlength{\parindent}{0pt}
\setlength{\parskip}{6pt}

\titleformat{\section}
{\normalfont\bfseries\fontsize{12}{14}\selectfont}
{}{0pt}{}
\titlespacing*{\section}{0pt}{6pt}{0pt}

\author{Jenny Sofía Morales}

\begin{document}

\begin{center}
    {\bfseries\fontsize{12}{14}\selectfont
        CÓMO GARANTIZAR LA CALIDAD EN UN PROYECTO DE TI
        MANTENER CALIDAD EN PROYECTOS DE TI
    }

    \vspace{0.4cm}

    {\itshape J. S. Morales López}

    \vspace{0.2cm}

    {\itshape 7690-08-6790 Universidad Mariano Gálvez}

    {\itshape Seminario de tecnologías de información}

    {\itshape
        \href{mailto:jmorales115@miumg.edu.gt}
        {jmorales115@miumg.edu.gt}
    }

    {\itshape
        \href{https://github.com/sofiaMorales3805/ForoAcademico-Seminario}
        {Repositorio en GitHub - Sofia Morales}
    }
\end{center}

\vspace{0.5cm}
\section*{Resumen}

Para garantizar la calidad en un proyecto de Tecnologías de la Información (TI)
implica planificar, diseñar, ejecutar y evaluar actividades que nos permitan cumplir los
requisitos del cliente y comprobar que la solución tecnológica presentada funciona de
manera adecuada. La calidad no depende únicamente de realizar pruebas al
final del desarrollo, también es establecer criterios de aceptación desde el
inicio, poder asignar responsabilidades, identificar riesgos y verificar los
resultados de forma continua. La gestión de la calidad comprende la
planificación, el aseguramiento de los procesos y el control de los
entregables. Para orientar estas actividades pueden emplearse referentes
como ISO 9001, dedicada a los sistemas de gestión de la calidad, e
ISO/IEC 25010, que ofrece un modelo para evaluar la calidad de productos
tecnológicos y de software. Asimismo, las prácticas de pruebas descritas
por ISTQB y las recomendaciones del NIST contribuyen a detectar defectos
y reducir riesgos de seguridad. Se concluye que la calidad se construye
durante todo el ciclo de vida del proyecto mediante objetivos medibles,
colaboración con los interesados, verificación continua y mejora de los
procesos, sin asumir que sea posible eliminar por completo los errores.

\noindent\textbf{Palabras clave:} calidad, proyectos de TI, aseguramiento de la calidad,
control de calidad, pruebas de software, gestión de riesgos.

\section*{Desarrollo del tema}

La calidad de un proyecto de TI puede entenderse como el grado en que
sus resultados satisfacen los requisitos acordados y las necesidades
reales de los usuarios. No es tan sensillo como solo entregar un sistema en el tiempo
establecido en el plan y el presupuesto: también debe ser útil, confiable, seguro y
adecuado. La gestión de proyectos requiere coordinar actividades, recursos, además de 
atender las dependencias tecnológicas y los riesgos que pueden afectar los entregables.\textsuperscript{[1,2]}

\noindent\textbf{Definición de requisitos y criterios de calidad}

El primer paso consiste en acordar con el cliente qué debe hacer la
solución y bajo qué condiciones debe funcionar. Los requisitos 
funcionales describen los servicios esperados; los no funcionales
establecen atributos como rendimiento, seguridad, disponibilidad y
facilidad de uso. Cada requisito importante debe contar con un criterio
de aceptación verificable. Por ejemplo, no es suficiente indicar que
una aplicación debe ser rápida: es necesario definir cómo se medirá su
tiempo de respuesta y en qué condiciones se evaluará. La norma
ISO/IEC 25010:2023 proporciona un modelo de calidad del producto que
puede servir como referencia para especificar y evaluar estas
características.\textsuperscript{[5]}

\noindent\textbf{Planificación de la calidad}

Una vez definidos los requisitos, el proyecto necesita un plan de
calidad que indique qué se revisará, quién será responsable, qué
métodos se aplicarán y qué evidencias se conservarán. Este plan puede
incluir revisiones de requisitos, inspecciones técnicas, pruebas,
criterios de aprobación y procedimientos para corregir defectos. El
Project Management Institute (PMI) presenta la planificación, el
aseguramiento y el control como componentes complementarios de la
gestión de la calidad. De este modo, el cumplimiento de los requisitos
se convierte en una responsabilidad organizada y no en una comprobación
improvisada al concluir el proyecto.\textsuperscript{[3]}

\noindent\textbf{Aseguramiento y control de calidad}

El aseguramiento de la calidad se orienta a prevenir problemas mediante
procesos adecuados, buenas prácticas y revisiones periódicas. En
cambio, el control de calidad evalúa los entregables para identificar
incumplimientos y defectos concretos. Por ejemplo, revisar los
procedimientos de desarrollo y la claridad de los requisitos ayuda al
aseguramiento; ejecutar casos de prueba y registrar los resultados
corresponde al control. Ambos enfoques son necesarios, porque un proceso
ordenado no elimina automáticamente los errores y una prueba final no
corrige las deficiencias de planificación.\textsuperscript{[3,4]}

\noindent\textbf{Pruebas y verificación continua}

Las pruebas deben incorporarse durante el ciclo de desarrollo y no
limitarse a la entrega final. Según el enfoque formativo de ISTQB, la
evaluación puede incluir revisiones estáticas, pruebas unitarias,
pruebas de integración, pruebas de sistema y pruebas de aceptación.
Cada nivel permite examinar riesgos distintos: desde el comportamiento
de un componente hasta la capacidad de la solución completa para
responder a las necesidades del usuario. Cuando sea conveniente, la
automatización de pruebas y su ejecución frecuente después de los
cambios ayudan a identificar regresiones, pero no reemplazan el
análisis y el criterio del equipo.\textsuperscript{[6]}

\noindent\textbf{Medición de resultados y seguimiento}

Para saber si se mantiene la calidad se deben utilizar indicadores
relacionados con los objetivos del proyecto. Entre ellos pueden
considerarse el porcentaje de pruebas aprobadas, la cantidad de
defectos pendientes según su severidad, el tiempo de resolución de
incidentes y el cumplimiento de criterios de aceptación. En proyectos
con servicios en operación también pueden evaluarse la disponibilidad
y los tiempos de respuesta. Las metas deben establecerse según el
contexto, la criticidad y los acuerdos con el cliente; una cifra
aislada no demuestra por sí sola que el producto tenga calidad.
\textsuperscript{[2,5]}

\noindent\textbf{Gestión de riesgos y seguridad}

La calidad incluye prevenir fallos que afecten la continuidad, los
datos y la confianza de los usuarios. Por ello, desde la planificación
se deben identificar dependencias, cambios de alcance, problemas de
integración y amenazas de seguridad, evaluando su probabilidad e
impacto. El marco Secure Software Development Framework (SSDF) del
NIST recomienda integrar prácticas de desarrollo seguro en el ciclo
de vida, proteger los componentes del software, reducir
vulnerabilidades y responder a las que se descubran. Estas medidas
fortalecen la calidad, aunque no garantizan la ausencia absoluta de
riesgos.\textsuperscript{[1,7]}

\noindent\textbf{Participación del equipo y mejora continua}

La comunicación entre desarrolladores, responsables del proyecto,
usuarios y especialistas de pruebas permite detectar interpretaciones
incorrectas antes de que se conviertan en problemas mayores. En
proyectos ágiles, la Guía de Scrum establece que la Definición de
Terminado debe expresar las medidas de calidad requeridas para
considerar completo un incremento del producto. Además, las
retrospectivas permiten revisar los resultados y mejorar los métodos
de trabajo. Aunque no todos los proyectos utilizan Scrum, la revisión
periódica de las evidencias y las lecciones aprendidas puede aplicarse
a distintos enfoques de gestión.\textsuperscript{[8]}

\section*{Observaciones y comentarios}

Considero que uno de los errores más frecuentes en los proyectos de TI
es asociar la calidad únicamente con la ausencia de errores visibles.
Un sistema puede superar varias pruebas y, aun así, presentar
problemas de seguridad, rendimiento o facilidad de uso que afecten a
las personas que dependen de él. Por esta razón, la evaluación debe
incluir tanto el funcionamiento técnico como la utilidad del producto
para el cliente.

También es importante reconocer que ninguna metodología garantiza por
sí sola el éxito de un proyecto. Las normas, los indicadores y las
herramientas aportan orientación, pero su utilidad depende de que el
equipo establezca objetivos realistas, registre evidencias y actúe a
tiempo ante los hallazgos. La prioridad debería ser prevenir defectos,
corregir sus causas y mejorar continuamente, en lugar de limitarse a
cumplir una lista de actividades.

\section*{Conclusiones}

\begin{enumerate}
    \item La calidad de un proyecto de TI debe definirse desde el inicio
    mediante requisitos claros, criterios de aceptación verificables y
    expectativas acordadas con el cliente.

    \item La planificación, el aseguramiento y el control de calidad son
    actividades complementarias: permiten prevenir problemas, supervisar
    los procesos y comprobar los resultados entregados.

    \item Las pruebas continuas y la medición de indicadores facilitan la
    identificación temprana de defectos y la toma de decisiones basadas
    en evidencias.

    \item La gestión de riesgos y la seguridad deben integrarse al
    proyecto, porque un producto funcional no necesariamente es seguro
    ni confiable en condiciones reales de operación.

    \item Garantizar la calidad no significa prometer cero errores, sino
    establecer un proceso sistemático de verificación, colaboración y
    mejora continua para alcanzar el nivel de calidad acordado.
\end{enumerate}

\section*{Referencias}

% Permite distribuir mejor los títulos largos y las direcciones URL.
\emergencystretch=1.5em
\begin{enumerate}[
    leftmargin=1.27cm,
    label={[\arabic*]},
    itemsep=6pt
]
    \item Carmona, R. (s. f.).
    \textit{5 elementos clave en la gestión de proyectos de TI}.
    MAGNET Telecommunications.
    \url{https://www.magnetmex.com/gestion-de-proyectos-de-ti/}

    \item Joseph, L. (s. f.).
    \textit{Gestión de proyectos de TI: guía completa para equipos y responsables}.
    Lemon Learning.
    \url{https://lemonlearning.com/es/blog/gestion-proyectos-ti-guia-completa}

    \item Rever, H. (2007, 14 de noviembre).
    \textit{Quality in project management---a practical look at chapter 8 of the PMBOK Guide}.
    Project Management Institute.
    \url{https://www.pmi.org/learning/library/quality-project-management-tools-7177}

    \item International Organization for Standardization. (2026).
    \textit{ISO 9001:2026: Quality management systems---Requirements}.
    \url{https://www.iso.org/standard/9001}

    \item International Organization for Standardization y
    International Electrotechnical Commission. (2023).
    \textit{ISO/IEC 25010:2023: Systems and software engineering---Systems
    and software Quality Requirements and Evaluation (SQuaRE)---Product quality model}.
    \url{https://www.iso.org/standard/78176.html}

    \item International Software Testing Qualifications Board. (s. f.).
    \textit{Certified Tester Foundation Level (CTFL) v4.0}.
    \url{https://istqb.org/certifications/certified-tester-foundation-level-ctfl-v4-0/}

    \item Souppaya, M., Scarfone, K., \& Dodson, D. (2022).
    \textit{Secure Software Development Framework (SSDF) Version 1.1:
    Recommendations for Mitigating the Risk of Software Vulnerabilities}
    (NIST SP 800-218).
    National Institute of Standards and Technology.
    \url{https://doi.org/10.6028/NIST.SP.800-218}

    \item Schwaber, K., \& Sutherland, J. (2020).
    \textit{The Scrum Guide: The definitive guide to Scrum---The rules of the game}.
    \url{https://scrumguides.org/scrum-guide.html}
\end{enumerate}

\end{document}
